\documentclass[10pt,a4paper]{amsart}
\usepackage[margin=19mm]{geometry}
\usepackage{amsmath,amssymb,mathtools}
\usepackage[scr=rsfs,cal=euler]{mathalfa}
\usepackage{xfrac}
\usepackage[unicode,hidelinks,bookmarks=false]{hyperref}
\newcommand{\ie}{i.e.\ }
\newcommand{\cf}{cf.\ }
\newcommand{\resp}{respectively\ }
\newcommand{\Subsection}{\subsection}
\providecommand{\nobreakdash}{\nobreak\hbox{-}}
\setlength{\emergencystretch}{2em}
\multlinegap0pt
\usepackage{amssymb}
\usepackage{xcolor}
\usepackage[normalem]{ulem}
\usepackage{bm}
\newtheorem{lem}{Lemma}
\newtheorem{thm}{Theorem}
\newtheorem{remark}{Remark}
\title{Symbolic Integration of Differential Forms:\\ from Abel to Zeilberger}
\author{Shaoshi Chen, David A. Cox, Yisen Wang}
\date{Source: arXiv:2601.00721v2 (CC BY 4.0)}
\begin{document}
\maketitle

\noindent\emph{Source and licence.} Symbolic Integration of Differential Forms:\\ from Abel to Zeilberger. Version arXiv:2601.00721v2 (CC BY 4.0), distributed by arXiv under the Creative Commons Attribution 4.0 International licence, http://creativecommons.org/licenses/by/4.0/, as recorded in the arXiv licence metadata for this exact version and shown on the abstract page https://arxiv.org/abs/2601.00721v2. Full licence notice: \texttt{licenses/arxiv-2601.00721-CC-BY-4.0.txt}.\par\smallskip

\noindent\textit{Editorial scope.} Counted unit: the article's thm thm (source0.tex lines 468--470). The article's complete original proof of it is delivered with it (source0.tex lines 472--517, one \texttt{proof} environment of 46 lines). No mathematics is written, repaired or re-derived: the statement and the proof are the article's own lines, in the article's own notation, and no retained line is substituted or re-flowed.  Retained uncounted as the source-local context the proof needs: lines 317--327; lines 347--353; lines 369--379; lines 407--424; lines 435--436.  Nothing was omitted from any retained span, so the package carries no omission record.  No retained line contains a citation, so the wrapper carries no bibliography and no bibliography entry.  No retained line contains a figure environment, an image-inclusion command or drawing code, so no image is delivered and the package declares no asset dependency.  Editorial formatting only, in addition: an AMS \texttt{amsart} preamble that restates the article's own declarations and macros used by the retained lines, namely the environments \texttt{lem}, \texttt{thm}, \texttt{remark} on separate counters, together with the standard packages those lines need.  The retained environments print in this document as Lemma 1, Theorem 1, Remark 1 in order of appearance, whereas the article prints the same items with the numbering of its own full document.  The complete native source of the article as distributed in the arXiv e-print for this exact version is bundled unchanged as \texttt{sources/192019/source0.tex}.
\begin{lem}
\label{LEM:form}
If $\omega = \omega_{m-1} + \omega^m  \in \bigwedge^p \mathcal{M}$ is as above, then we have
\begin{itemize}
\item[{\rm(i)}] $d^m (\omega^m) = 0$.
\item[{\rm(ii)}] $d(\omega \,dx_m) = d_{m-1}(\omega)\, dx_m$.
\item[{\rm(iii)}] $d\hskip1pt\omega =0$ if and only if  $d_{m-1}(\omega_{m-1}) = 0$ and $d^m(\omega_{m-1})  + d_{m-1}(\omega^m)=0$.
\item[{\rm(iv)}] If $\omega$ is free of $dx_m$, then $d\hskip1pt\omega = 0$ if and only if  $d_{m-1}(\omega) = 0$ and $d^m (\omega) = 0$.
\item[{\rm(v)}] If $\omega $ is free of $x_m$ and $dx_m$, then $\omega = d\mu$ for some $(p-1)$-form $\mu$ if and only if $ \omega = d_{m-1}(\tilde \mu)$ for some $(p-1)$-form $\tilde \mu$ also free of $x_m$ and $dx_m$.
\end{itemize}
\end{lem}

\begin{equation}\label{EQ:partialfm} 
\begin{aligned}
 \partial_m(f_i) &= \partial_i(f_m) = \partial_i(\partial_m(g_m))= \partial_m(\partial_i(g_m))\\
  &= \partial_m\bigg(\partial_i(a_m) + \sum_j \Big(c_{m, j}\frac{\partial_i(b_{m, j})}{b_{m, j}} + \partial_i(c_{m, j}) \log b_{m, j}\Big)\!\bigg) \\
 &= \partial_m\bigg(\partial_i(a_m) + \sum_j c_{m, j} \frac{\partial_i(b_{m, j})}{b_{m, j}} \bigg) + \sum_j \partial_i(c_{m, j})\frac{\partial_m(b_{m, j})}{b_{m, j}},
\end{aligned}
\end{equation}

\begin{thm}
\label{THM:1form}
Let $\omega = f_1\, dx_1 + \cdots + f_m \, dx_m$ be a closed $1$-form with $f_i \in K = k(x_1,\dots,x_m)$. 
%{\color{red}
Then $\omega = dg$ for some $g$ of the form
\[
g = a + \sum_{i=1}^m \sum_j c_{i,j} \log b_{i,j},
\]
where  $a \in K$, $c_{i,j} \in  \bar k$, and $b_{i,j} \in k(x_1,\ldots,x_{i-1},c_{i,j})[x_i]$ has positive degree in~$x_i$.
%are pairwise coprime polynomials of positive degree.
\end{thm}

\begin{remark}\label{RMK:proper} 
By convention, we define~$\deg_{x_i}(0) := - \infty$ for each~$i=1,\ldots,m$. Hence $0$ is proper in~$x_i$.  %When we regard $K$ as a vector space over the subfield $K_i  := k(x_1, \ldots, x_{i-1}, x_{i+1}, \ldots, x_m)$, the set of all proper rational functions in $x_i$ forms a subspace of $K$ over $K_i$.
Also,
\begin{itemize}
\item For any nonzero $f \in {\bar k}(x_1,\ldots,x_m)$, $\partial_i(f)/f$ is proper in $x_i$ and has a squarefree denominator in ${\bar k}[x_1,\dots,x_m]$.
\item For any monic $f \in {\bar k}(x_1,\ldots,x_{m-1})[x_m]$, $\partial_i(f)/f$ is proper in $x_m$.
%The field $K$ can be viewed as a vector space over its subfield 
%\[K_i := k(x_1, \ldots, x_{i-1}, x_{i+1}, \ldots, x_m),\] 
%which is the constant field of $K$ with respect to the derivation $\partial_i$.
%The set of all proper rational functions in $x_i$ forms a subspace of $K$ over $K_i$.
%Any rational function $f\in K$ can uniquely written as $f = p + a/b$ such that $p \in K_i[x_i]$
%and $a/b$ is proper in~$x_i$. We call $p$ and $a/b$ the \textit{polynomial part} and \textit{proper part} of $f$ in $x_i$, respectively. 
\item For any proper $a/b\in {\bar k}(x_1,\ldots.x_m)$ in $x_i$, we have 
\[\partial_j\Big(\frac{a}{b}\Big) =\frac{\partial_j(a)\hskip.5pt b - a\hskip.5pt \partial_j(b)}{b^2}. \]
So $\partial_j(a/b)$ is still proper in $x_i$. %This implies that the operations of taking the polynomial part and the proper part in $x_i$ commute with any derivation $\partial_j$ with $j\in \{1, \ldots, m\}$. If $f$ is proper in $x_i$ and $f = \partial_i(g)$ for some $g\in K$, then there exists a unique proper $h\in K$ in $x_i$  such that $f = \partial_i(h)$. Indeed, $h$ can be taken as the proper part of $g$. Consequently,  
\item For any $f\in {\bar k}(x_1,\ldots,x_m)$,  $f =0$ if and only if $f$ is proper in $x_i$ and lies in ${\bar k}(x_1,\ldots,x_{i-1},x_{i+1},\ldots,x_m)$.
\end{itemize}
\end{remark}

\begin{remark}\label{RMK:sumsimple}  Note that the sum of two simple forms is again a simple form since proper rational functions are closed under addition and the least common multiple of squarefree polynomials in $k[x_1,\ldots,x_m]$ is still squarefree. 
\end{remark}

\begin{thm}
Algorithm {\tt HermiteOneForm} is correct.
\end{thm}

\begin{proof} We focus on the case $m > 1$.  In Step (2), we write $f_m = \partial_m(a_m) + \bar f_m$ with $a_m, \bar{f}_m \in K$ and $\bar{f}_m$  proper in $x_m$, and in Step (3), $\bar{f}_m = \sum_j c_{m, j} \frac{\partial_m(b_{m, j})}{b_{m,j}} = \partial_m\big(\sum_j c_{m, j} \log b_{m, j}\big)$.
The discussion following \eqref{EQ:partialfm} shows that the $c_{m,j}$ lie in $\bar k$.  

Step (4) defines $\bar f_i := \sum_j c_{m, j} \frac{\partial_i(b_{m, j})}{b_{m,j}} = \partial_i\big(\sum_j c_{m, j} \log b_{m, j}\big)$ for $i = 1,\ldots,m-1$.  Note that $\bar f_i \in K$ by the proof of Theorem~\ref{THM:1form}. Furthermore,
\begin{equation}\label{EQ:fifm}
\partial_m(\bar f_i) = \partial_m\partial_i\big({\textstyle\sum_j} c_{m, j} \log b_{m, j}\big) = \partial_i\partial_m\big({\textstyle\sum_j} c_{m, j} \log b_{m, j} \big) = \partial_i(\bar f_m).
\end{equation}

Now consider the 1-form $\bar{r}_m := \bar{f}_1 \, dx_1 + \cdots + \bar{f}_m \, dx_m$
defined in Step (4). We claim that $\bar r_m$ is a closed form.  It suffices to prove that  
\[
    \partial_i\bigl(\bar{f}_j\bigr) = \partial_j\bigl(\bar{f}_i\bigr), \ \text{for any}\ 1 \leq i < j \leq m-1.
\]
Since $\partial_m\bigl(\bar{f}_i\bigr) = \partial_i\bigl(\bar{f}_m\bigr)$ for $i = 1,\ldots,m-1$ by \eqref{EQ:fifm}, we have
\[
    \partial_m \bigl(\partial_i(\bar{f}_j)-\partial_j(\bar{f}_i)\bigr) = \partial_i(\partial_j(\bar{f}_m))-\partial_j(\partial_i(\bar{f}_m)) = 0.
\]
By Remark~\ref{RMK:proper}, $\partial_i(\bar{f}_j)-\partial_j(\bar{f}_i)$ is proper in~$x_m$ (each $b_{m,j}$ is monic in $x_m$), and hence must be zero. So the claim holds.   Remark~\ref{RMK:proper} also implies $\bar f_i$ is proper in $x_i$ with squarefree denominator for $i = 1,\dots,m$.  It follows that $\bar{r}_m$ is a simple form with coefficients in $K$.

%To complete the proof that $\bar r_m$ is simple, we need to show that 
%Let $e_{m,j}$ denote the leading coefficient of $b_{m,j}$ with respect to~$x_m$. Then,
%the proper part of the logarithmic derivative ${\partial_i(b_{m,j})}/{b_{m,j}}$ in~$x_m$ is easily seen to be 
%\[
%    \frac{\partial_i(b_{m,j})}{b_{m,j}} - \frac{\partial_i(e_{m,j})}{e_{m,j}}.
%\]
% Since $e_{m,j} \in k[x_1,\ldots,x_{m-1}]$ is free of~$x_m$, it holds that   
% \[
%   \partial_m\biggl( \frac{\partial_i(b_{m,j})}{b_{m,j}} - \frac{\partial_i(e_{m,j})}{e_{m,j}} \biggr) =  \partial_m\biggl( \frac{\partial_i(b_{m,j})}{b_{m,j}} \biggr) = \partial_i \biggl( \frac{\partial_m(b_{m,j})}{b_{m,j}}\biggr).
% \]
% By the linearity of differentiation,
% \[
%   \partial_m\Biggl(\sum_j c_{m,j} \biggl(\frac{\partial_i(b_{m,j})}{b_{m,j}} - \frac{\partial_i(e_{m,j})}{e_{m,j}} \biggr) \Biggr) = \partial_i \Biggl( \sum_j c_{m,j}\frac{\partial_m(b_{m,j})}{b_{m,j}}\Biggr) = \partial_i\bigl(\bar{f}_m \bigr).
% \]
%By the linearity of the operation of taking the proper part, we  have 
%\[
%    \bar{f}_i = \sum_{j} c_{m,j} \biggl(\frac{\partial_i(b_{m,j})}{b_{m,j}} - \frac{\partial_i(e_{m,j})}{e_{m,j}}\biggr),
%\]
%which is proper in~$x_i$ with a squarefree denominator in~$k[x_1,\ldots,x_m]$.
%Therefore, the form $\bar{r}_m$ constructed in Step~$(3)$ is simple. 

The form $\omega_{m-1}:=\omega-da_m-\bar{r}_m$ in Step~(4) is closed since~$d a_m$ is exact and both $\omega$ and $\bar{r}_m$ are closed. (Note that $\omega_{m-1}$ is the form $\widetilde \omega$ used in the proof of Theorem~\ref{THM:1form}.) Furthermore, $\omega_{m-1}$ involves no $dx_m$ terms and then is free of~$x_m$ by {\rm(iv)} of Lemma~\ref{LEM:form}.  Thus $\omega_{m-1}$ is a closed 1-form over $F = k(x_1,\dots,x_{m-1})$.

In Step~(5), recursively applying the above procedure over~$F$ leads to $\omega_{m-1}= dg_{m-1}+r_{m-1}$,
where $g_{m-1}\in F$ and $r_{m-1}$ is a simple form over $F$.
Set $g := a_m+g_{m-1}$ and $r := \bar{r}_m + r_{m-1}$. Thus $\omega = dg + r$, where $r$ is simple by Remark~\ref{RMK:sumsimple}. 
\end{proof}

\end{document}
